% !PRINTED 92
\weberpage{92}
%<<CONT

Bedingung eine neue der Grössen $y$ enthält, von einander unabhängig sind. \gesperrt{Ist eine dieser Bedingungen nicht erfüllt, so hat das gegebene Gleichungssystem keine Lösung.}

Wenn wir in dem Gleichungssystem (1) die $x$ nicht als Unbekannte, sondern als Veränderliche betrachten, so werden auch die $y$ veränderliche Grössen sein. Bei dieser Auffassung nennen wir das System (1) eine \gesperrt{lineare Substitution}, insofern dadurch der Uebergang von einem System von Variablen zum anderen vermittelt wird. Nur wenn die Determinante $A$, die wir jetzt die \gesperrt{Substitutions-Determinante} nennen, von Null verschieden ist, werden die $y$ als unabhängige Variable angesehen werden können. Ist dies der Fall, so ergiebt das System (5) die Darstellung der Variablen $x$ durch die $y$ oder die zu (1) \gesperrt{inverse Substitution}.


\begin{center}
{\large \S.\ 27.}

\vspace{0.5em}
{\large\gesperrt{Multiplication von Determinanten.}}
\end{center}

Der Satz, den wir jetzt noch beweisen wollen, lehrt, wie man das Product zweier Determinanten von gleich viel Reihen durch eine einzige Determinante von ebenso viel Reihen darstellen kann. Man wird am einfachsten darauf geführt, wenn man die Auflösung von zwei Systemen linearer Gleichungen betrachtet.

Es seien jetzt die Coëfficienten $a_i^{(k)}, b_i^{(k)}$, wenn $i,k$ die Reihe der Zahlen $1,2,\ldots,n$ durchlaufen, beliebige \gesperrt{veränderliche Grössen}, ebenso die Grössen $x_i,y_i,z_i$, die nur an die Relationen gebunden sind:
\begin{align}
 \sum_i a_i^{(k)}x_i&=y_k,\tag{1}\\
 \sum_k b_h^{(k)}y_k&=z_h.\tag{2}
\end{align}
Wenn nun die Aufgabe gestellt wird, die Variablen $x$ durch die Variablen $z$ zu bestimmen, so kann diese Aufgabe auf doppelte Art gelöst werden. Man kann nach §. 26, (4) die Gleichungen (1) in Bezug auf $x$, die Gleichungen (2) in Bezug auf $y$ auflösen, und die letzteren Ausdrücke in die ersteren einsetzen. Bezeichnen wir mit $A,B$ die Determinanten der beiden Gleichungssysteme, also:
%CONT>>
